\documentclass[a4paper,12pt,openright,titlepage,oneside]{book}

% \usepackage[english,brazil]{babel}
% Define regra de gramática para separar sílabas {babel}
% e altera os títulos (como chapters, sections, references) para português
\usepackage[brazil]{babel}

% Define uso de caracteres acentuados no PDF gerado
% e permite copiar corretamente o texto do PDF.
\usepackage[T1]{fontenc}
\usepackage[utf8]{inputenc}

% Adição do pacote do template da UnB
\usepackage{template-FT-UnB/ft2unb}

% Define a fonte Times New Roman para o documento
\usepackage{mathptmx}

\DeclareGraphicsExtensions{.jpg,.pdf,.mps,.png,.gif,.eps}
\graphicspath{imagens} % Define diretório de imagens

% Arquivo com lista de hifenização correta de algumas palavras.
% Defina novas palavras no arquivo à medida que verificar que a hifenização automática
% está errada para tais palavras.
\input{hifenizacao}

% ALTERE OS VALORES DENTRO DAS CHAVES DOS COMANDOS NESTA SEÇÃO PARA INCLUIR OS SEUS DADOS E DADOS DA SUA
% DISSERTAÇÃO DE MESTRADO OU TESE DE DOUTORADO
% -----------------------------------------------------------------------------------------------------

% \onehalfspacing
\title{TÍTULO DA DISSERTAÇÃO/TESE}
\author{NOME DO AUTOR}
\date{2011-03-16} % Data da defesa

\grau{Mestre} % Mestre ou Doutor
\area{Transportes} % Nome do curso
\siglaarea{ENC} % Sigla do departamento (ENC = Engenharia Civil e Ambiental)
\tipodemonografia{Dissertação} % Dissertação ou Tese
\programa{Mestrado} % Mestrado ou Doutorado
\autorendereco{ENDEREÇO DO AUTOR.} % Endereço do autor da dissertação/tese
\totalpgs{147} % Total de páginas atualmente na sua dissertação
\dia{16} % Dia da defesa
\mes{Junho} % Mês da defesa
\ano{2011} % Ano da defesa
\numpublicacao{T.DM-XXX/AAAA} % Número da publicação (após a defesa, tal número deve ser obtido na secretaria)

\titulolinhai{TÍTULO DA DISSERTAÇÃO/TESE}
\titulolinhaii{}
\titulolinhaiii{}
\titulolinhaiv{}

\autori{NOME DO AUTOR}
% Caso seu nome não caiba em uma única linha, divida-o nos comandos abaixo
% \autorii{}
% \autoriii{}

\membrodabancai{Prof. Dr. NOME DO ORIENTADOR, Titulação (Instituição)}
\membrodabancaifuncao{Orientador}
\membrodabancaii{Prof. Dr. NOME DO EXAMINADOR, Titulação (Instituição)}
\membrodabancaiifuncao{Examinador interno}
\membrodabancaiii{Prof. Dr. NOME DO EXAMINADOR, Titulação (Instituição)}
\membrodabancaiiifuncao{Examinador externo}
\membrodabancaiv{}
\membrodabancaivfuncao{}
\membrodabancav{}
\membrodabancavfuncao{}
% -----------------------------------------------------------------------------------------------------

% line-numbers, inputencoding=utf8/latin1
% Define o estilo para listagens de código fonte
\lstset{
  numbers=left, % Numeração de linhas à esquerda
  stepnumber=1,
  firstnumber=1,
  numberstyle=\tiny,
  extendedchars=true,
  frame=none,
  basicstyle=\footnotesize,
  stringstyle=\ttfamily,
  showstringspaces=false,
  % language=Java, % Deve ser definida na inclusão de cada trecho de código, pois podem existir linguagens diferentes em exemplos diferentes
  breaklines=true,
  breakautoindent=true,
  % Estilos de comentário de uma e várias linhas
  morecomment=[l]{--}, morecomment=[s]{/*}{*/}, morecomment=[s]{<!--}{-->}, morecomment=[s]{--[[}{--]]}
}

% Adição de metadados no PDF (propriedades do documento PDF)
\makeatletter
\hypersetup{
  pdftoolbar=true,        % Show Acrobat’s toolbar?
  pdfmenubar=true,        % Show Acrobat’s menu?
  pdffitwindow=false,     % Window fit to page when opened
  pdfstartview={FitH},    % Fits the width of the page to the window
  pdftitle={\@title},
  pdfauthor={\@author},
  pdfsubject={\tipodemonografianome\ de\ \programastr\ em\ \areastr}, % Subject of the document
  pdfcreationdate={\pdfdate}
}
\makeatother

\makeindex
\makenomenclature % Necessário para gerar lista de siglas

\begin{document}

  % Ordem dos elementos pré-textuais conforme PPGT:
  % DEDICATÓRIA (opcional)
  \include{capitulos/agradecimentos}
  % RESUMO
  \include{capitulos/resumo}
  % ABSTRACT
  \include{capitulos/abstract}

  % ÍNDICE (sumário)
  \pdfbookmark[0]{Sumário}{sumario}
  \sumario

  % LISTA DE TABELAS
  \pdfbookmark[0]{Lista de Tabelas}{listatabelas}
  \listadetabelas

  % LISTA DE FIGURAS
  \pdfbookmark[0]{Lista de Figuras}{listafiguras}
  \listadefiguras

  % LISTA DE QUADROS (PPGT)
  \pdfbookmark[0]{Lista de Quadros}{listaquadros}
  \listadequadros

  % LISTA DE SÍMBOLOS, NOMENCLATURA E ABREVIAÇÕES
  \renewcommand{\nomname}{LISTA DE SÍMBOLOS, NOMENCLATURA E ABREVIAÇÕES} % Define um caption à lista de siglas
  \pdfbookmark[0]{Lista de Termos e Siglas}{nomenclatura}
  \printnomenclature[2.5cm]

  \mainmatter % Inicia a numeração normal cardinal
  \setcounter{page}{1} \pagenumbering{arabic}
  % Numeração de páginas no rodapé, alinhada à direita (PPGT 2023)
  \pagestyle{fancy}
  \fancyhf{}
  \renewcommand{\headrulewidth}{0pt}
  \fancyfoot[R]{\thepage}
  \fancypagestyle{plain}{%
    \fancyhf{}%
    \renewcommand{\headrulewidth}{0pt}%
    \fancyfoot[R]{\thepage}%
  }

  \include{capitulos/introducao}
  \chapter{TRÂNSITO E SEGURANÇA VIÁRIA} \label{capitulo1}

\section{ACIDENTES DE TRÂNSITO E GESTÃO DA VELOCIDADE}

Descreva aqui o contexto sobre acidentes de trânsito e a importância da gestão de velocidade.

\section{FISCALIZAÇÃO ELETRÔNICA DE VELOCIDADE}

Descreva os sistemas de fiscalização eletrônica e seu papel na segurança viária.

\begin{quadro}[htbp]
\caption{Principais estudos sobre efeitos da fiscalização de velocidade na segurança viária, no âmbito internacional}
\label{quad:estudos_internacionais}
\centering
\small
\begin{tabular}{|p{4cm}|p{5cm}|p{4cm}|}
\hline
\textbf{Autor} & \textbf{Método} & \textbf{Resultados} \\
\hline
Autor 1 (Ano) & Descrição do método & Principais achados \\
\hline
Autor 2 (Ano) & Descrição do método & Principais achados \\
\hline
\end{tabular}
\end{quadro}

\section{TÓPICOS CONCLUSIVOS}

Apresente os pontos conclusivos do capítulo.

  \chapter{MATERIAIS E MÉTODOS}
\label{capitulo2}

\section{MATERIAIS}

Descreva os materiais utilizados na pesquisa.

\section{MÉTODOS}

Descreva a metodologia geral da pesquisa.

\subsection{Metodologia DNER (1998)}

Descreva a metodologia DNER utilizada.

\begin{table}[htbp]
  \centering
  \caption{Índices para soma ponderada}
  \label{tab:indices_ponderada}
  \small
  \begin{tabular}{lc}
    \hline
    \textbf{Tipo de Acidente} & \textbf{Índice} \\
    \hline
    Com vítima fatal & 13 \\
    Com ferido & 3 \\
    Sem vítima & 1 \\
    \hline
  \end{tabular}
  \\[0.2cm]
  {\small Fonte: DNER (1998).}
\end{table}

\subsection{Metodologia Vadeby e Forsman (2018)}

Descreva a metodologia Vadeby e Forsman.

\subsubsection{Subtítulo do 3º nível}

O 3º nível pode ser itemizado por letras (a, b, c...) ou por números romanos (i, ii, iii...). Uma vez escolhida a forma de identificação dos itens, este deve ser mantido para todo o texto. Apenas o primeiro e o segundo níveis são numerados no sumário.

\subsection{Teste t de Student}

Descreva o teste estatístico utilizado.

  \chapter{ANÁLISE DOS RESULTADOS}
\label{capitulo3}

\section{DELINEAMENTO DO CENÁRIO DE ANÁLISE}

Descreva o cenário de análise.

\subsection{Descrição das bases de dados}

Descreva as bases de dados utilizadas.

\subsection{Levantamento de dados}

Descreva como foi feito o levantamento de dados.

\paragraph{(a) Dados tipo I}

Descreva o tipo de dado coletado.

\section{APLICAÇÃO DOS MÉTODOS}

Apresente a aplicação dos métodos descritos no capítulo anterior.

\subsection{Efeito na severidade}

Analise o efeito na severidade dos acidentes.

\subsection{Relação dos resultados com as características dos locais}

Relacione os resultados obtidos com as características dos locais estudados.

\subsection{Comparação com outros estudos}

Compare os resultados com outros estudos da literatura.

  \include{capitulos/conclusao}

  % \bibliographystyle{bibliografia/abnt-num} % Estilo bibliográfico ABNT numérico
  \bibliographystyle{bibliografia/abnt-alf} % Estilo bibliográfico ABNT alfabético (PPGT)
  % \bibliographystyle{bibliografia/sbc} % Estilo bibliográfico da Sociedade Brasileira de Computação (SBC)

  % \renewcommand{\bibname}{REFERÊNCIAS BIBLIOGRÁFICAS} % Define o caption da seção de bibliografia
  % \addcontentsline{toc}{chapter}{REFERÊNCIAS BIBLIOGRÁFICAS}

  % Não pode ter espaço entre os nomes dos arquivos bib
  \bibliography{bibliografia/referencias}

  \include{capitulos/apendice}
\end{document}
